\section{Introduction}\label{sec:introduction}

A lower bound for scalar curvature is a local condition. Gromov's
macroscopic dimension conjecture predicts a global consequence: at a
controlled scale, two dimensions can be collapsed. We prove the
quantitative form of this assertion for arbitrary complete manifolds
in dimensions at least four.

Throughout, scalar curvature means the ordinary pointwise Riemannian
invariant. Fixed-scale macroscopic scalar curvature uses ball-volume
comparison at a prescribed positive radius
\citep[Definitions~1.2--1.3]{kumarSen2026Width}.

For a metric space \(X\) and an integer \(d\ge0\), define its
\emph{Urysohn \(d\)-width} by
\[
 \UW_d(X)=\inf_f\ \sup_{y\in K}\diam_X f^{-1}(y),
\]
where the infimum runs over continuous maps \(f:X\to K\) to simplicial
complexes of dimension at most \(d\). The complexes have their usual
polyhedral topology. The diameter is that of the entire fiber, including
all its components; empty fibers may be omitted.

\begin{theorem}\label{thm:main}
For every integer \(n\ge4\), there is a finite constant \(C_n\) such
that every connected complete smooth Riemannian \(n\)-manifold
\((M,g)\), without boundary and with \(\Scal_g\ge1\), admits a
continuous map \(f:M\to K\) to a simplicial complex of dimension at
most \(n-2\) satisfying
\[
 \diam_g f^{-1}(y)\le C_n\qquad(y\in K).
\]
In particular, \(\UW_{n-2}(M,g)\le C_n\).
\end{theorem}

The constant is independent of \(M\), its topology, and its geometry.
The theorem applies to compact and noncompact manifolds, with no
orientability, spin, covering, or bounded geometry assumption. If
\(\Scal_g\ge\sigma^2>0\), applying the theorem to \(\sigma^2g\) gives
\(\UW_{n-2}(M,g)\le C_n/\sigma\).

For a closed connected Riemannian manifold, let
\(\operatorname{FillRad}(M,g)\) be Gromov's filling radius
\citep{gromov1983Filling}, defined by the Kuratowski embedding and the
fundamental class, with coefficients in
\(\Z\) if \(M\) is orientable and in \(\Z_2\) otherwise
\citep[Definition~3.1]{nabutovskyRotmanSabourau2021}.
The next corollary gives the closed-manifold case in dimensions at least
four of Gromov's filling-radius Conjecture~\(A^{-}\)
\citep[\S3]{gromov2017Dozen}.

\begin{corollary}[Closed-manifold filling-radius bound]
\label{cor:filling-radius}
Let \((M,g)\) be a closed connected smooth Riemannian \(n\)-manifold
with \(n\ge4\) and \(\Scal_g\ge\sigma^2\) for some \(\sigma>0\). Then
\[
 \operatorname{FillRad}(M,g)\le \frac{C_n}{2\sigma},
\]
where \(C_n\) is the constant in Theorem~\ref{thm:main}.
\end{corollary}

\begin{proof}
Gromov's comparison, in the form recorded by Nabutovsky, Rotman, and
Sabourau \citep[Definition~6.1 and equation~(6.1)]{nabutovskyRotmanSabourau2021},
uses the same continuous-map and entire-fiber width convention as above.
A target of dimension at most \(n-2\) is also admissible for
\(\UW_{n-1}\). Since \(M\) is closed, the comparison and the scaled
form of Theorem~\ref{thm:main} give
\[
 \operatorname{FillRad}(M,g)
 \le \tfrac12\UW_{n-1}(M,g)
 \le \tfrac12\UW_{n-2}(M,g)
 \le \frac{C_n}{2\sigma}.
\]
\end{proof}

For the closed-cover consequence we also use the \emph{continuous
macroscopic dimension}: $\dim_{\mathrm{mc}}X\le d$ means that there is a
continuous map to a simplicial complex of dimension at most $d$, with its uniform simplex metric
and a finite uniform bound on the diameters of its entire fibers.
Each simplex has the same standard Euclidean shape and edge length.

\begin{corollary}[Continuous macroscopic dimension of universal covers]
\label{cover:main}
Let $(M,h)$ be a closed connected smooth Riemannian manifold of
dimension $n\ge2$ with positive scalar curvature.  Its universal cover,
with the lifted original metric $\widetilde h$, satisfies
\[
 \dim_{\mathrm{mc}}(\widetilde M,\widetilde h)\le n-2.
\]
\end{corollary}

\begin{corollary}[Aspherical obstruction]\label{cor:aspherical}
No closed smooth aspherical manifold of dimension \(n\ge4\) admits
a Riemannian metric of strictly positive scalar curvature.
\end{corollary}

\begin{proof}
It suffices to consider a connected manifold; passing to its orientation
double cover if necessary preserves the hypotheses. Suppose such a metric
exists. Its scalar curvature has a positive minimum, so the lifted metric
on the contractible universal cover \(X\) is complete with a uniform
positive scalar-curvature lower bound. The deck group acts isometrically
and cocompactly. The scaled form of Theorem~\ref{thm:main} gives a
continuous map from \(X\) to an \((n-2)\)-complex with uniformly bounded
entire fibers. The width and filling-radius implications in Gromov's
noncompact argument \(A^+\Rightarrow B\)
\citep[\S3, pp.~6--8]{gromov2017Dozen} exclude such a map on a
contractible universal cover with a cocompact deck action. This gives
the contradiction.
\end{proof}

An independent proof in every dimension follows from the stronger
rational-inessentiality theorem of
\citet[Theorem~1.1 and Corollary~1.2]{openai2026RationalInessentiality}.

\subsection{History and significance}

Gromov introduced macroscopic dimension in his study of positive scalar
curvature \citep[\S2]{gromov1996}. His quantitative conjecture already asks
whether a complete manifold with a positive scalar-curvature lower bound
can be mapped continuously to a polyhedron of codimension two with all
point inverses bounded at the curvature scale
\citep[\S\(2\frac12\)]{gromov1996}; see also
\citep[p.~3180, Comments after Conjecture~3.2]{gromov2019Mean}.
The diameter condition concerns each entire fiber, including all its
components. In the universal cover of a closed manifold the corresponding
macroscopic-dimension assertion asks only for some finite bound.
The continuous invariant and the version requiring a Lipschitz map to a
uniformly metrized complex are distinct
\citep[Definitions~1.1--1.2]{dranishnikov2014}.

The pointwise curvature hypothesis cannot be replaced by its fixed-scale
macroscopic counterpart. Kumar and Sen construct closed manifolds in
every dimension at least four with a prescribed positive fixed-scale
macroscopic scalar-curvature lower bound and arbitrarily large
Urysohn codimension-two width \citep[Theorem~A]{kumarSen2026Width}.

The low-dimensional cases have a substantial earlier history.
Gromov and Lawson obtained geometric width bounds for closed three-manifolds
\citep[Corollary~10.11]{gromovLawson1983}, and Bolotov proved that the
universal cover of a closed three-manifold cannot have macroscopic
dimension two \citep{bolotov2003}.
For closed three-manifolds with positive scalar curvature, Liokumovich
and Maximo constructed Morse functions with curvature-scale bounds for
the areas and ambient diameters of connected level components
\citep[Theorem~1.1]{liokumovichMaximo2023}.
Liokumovich and Wang constructed proper Morse functions on
complete noncompact three-manifolds with uniformly positive scalar
curvature, with uniform diameter bounds for connected level components
\citep[Theorem~1.1]{liokumovichWang2023}.
We use that result and an explicit passage to the Reeb graph for the
three-dimensional closed-cover consequence; the surface case follows
from the Bonnet--Myers argument.

The aspherical obstruction in dimensions four and five was established
by Chodosh and Li \citep[Theorem~2]{chodoshLi}, with an independent
proof in dimension five by Gromov
\citep[\S7, Main Corollary]{gromov2020Aspherical}.

For the macroscopic-dimension conjecture, Bolotov proved the closed spin
case under the Strong Novikov Conjecture and a cohomological-dimension condition
\citep{bolotov2009}. Bolotov and Dranishnikov then obtained the
codimension-two conclusion for closed spin manifolds under the
Strong Novikov Conjecture and
injectivity of the connective-to-periodic real K-homology map
\citep[Theorem~4.3]{bolotovDranishnikov2010}.
For closed orientable totally non-spin manifolds of dimension at least
five, they also showed that a macroscopic dimension drop by one implies
a drop by two \citep[Theorem~3.4]{bolotovDranishnikov2016}.
Daher and Dranishnikov established the analogous implication for closed
orientable totally non-spin four-manifolds with residually finite
fundamental group
\citep[Theorem~3.3]{daherDranishnikov2019}.
These topological advances concern universal covers and finite bounds;
the present quantitative statement is formulated for arbitrary complete
manifolds with a bound depending only on dimension after scalar normalization.


The geometric method uses the minimal-hypersurface approach of Schoen
and Yau, in which stability and the Gauss equation pass scalar-curvature
information to hypersurfaces \citep[Section~1]{schoenYau1979Structure},
and the torical symmetrization mechanism described by Gromov
\citep[Section~2]{gromov2018Metric}. Gromov's \(\mu\)-bubble method
uses prescribed mean curvature to localize the hypersurfaces entering
scalar-curvature descent \citep[\S\(5\frac56\)]{gromov1996}.
The weighted variational formulation and its relation to warped-circle
geometry and stability are developed by Chodosh and Li
\citep[Section~3 and Remark~11]{chodoshLi}.

Removing singular regions while retaining scalar-curvature information
is also central to Lohkamp's splitting theory
\citep[\S1.2 and Theorem~3.15]{lohkamp2022SplittingsII}.
In the construction below, profile updates exclude singular points from
the regions where the next scalar bound is required; the wall and both
retained sides are then regularized together.

Recent results also address quantitative width and singular descent.
Antonelli proves width bounds under a positive scalar-curvature lower
bound, an almost-nonnegative Ricci lower bound, and uniform noncollapse,
with constants depending on the volume lower bound
\citep[Theorem~1.4]{antonelli2026Volume}.
Bi and Zhu develop singular scalar-curvature descent for enlargeability
obstructions and cubical-width inequalities
\citep[Theorem~B and Section~4]{biZhu2026}.
The present construction retains both bulk sides and the wall through
countably many cuts, with compatible product ends.

\subsection{The proof and its main constructions}

Fix an original length \(D\), depending only on \(n\), and a separated
net \((x_i)\) in \((M,g)\). The inner balls of radius \(D\) cover \(M\),
and the closed outer balls of radius \(2D\) are locally finite. Each
center has its own label. We construct a continuous collapse map
$q:M\to M$ with $d_g(q(x),x)<1$, together with compatible walls
separating the inverse images of each inner ball and the exterior of
its outer ball. The map $q$ records original positions while cylinders
are inserted along the cuts. The final
walls have no common intersection of \(n-1\) distinct labels. Continuous
sign coordinates $t_i:M\to[0,1]$ are chosen to equal $1$ on the
$q$-preimage of the inner ball and $0$ on the $q$-preimage of the exterior
of its outer ball. Their transition regions lie in wall neighborhoods
chosen simultaneously to preserve the intersection bound. They give the
desired map: at most \(n-2\) coordinates
are fractional at a point, and some coordinate is one. Equality of all
coordinates confines both collapse images to the same original outer
ball, giving a bound for the entire fiber.

Three points are essential to make this construction uniform and to
allow arbitrary complete metrics.

First, the separations must become as wide as prescribed without
spending curvature in proportion to the number of labels. For a metric
\(h\) and a smooth function \(t\), put
\begin{equation}\label{intro:drift-scalar}
 \mathcal D(h,t)=\Scal_h+2\Delta_h t-|\dd t|_h^2.
\end{equation}
Fix a background metric \(h_0\) throughout a round. The local stretching
estimate changes \(h\) monotonically and increases \(t\), while paying
for its loss with a decrease in \(\tr_h h_0\).
The pressure law in \Cref{pressure:law} provides a quadratic budget in
\(t\). This makes the initial scale depend only on \(n\), even when
there is no bound on the overlap multiplicity of the original balls.
Trace losses telescope over a countable locally finite round. Each
label stops after its own finite number of rounds, so the entire
construction stabilizes on every compact set
(\Cref{rounds:wide}). Diagonal circle warps convert the drift scalar
bound to an ordinary scalar-curvature inequality on every required
finite product.

Second, a cut produces both extended bulk sides and a wall of one lower
base dimension. We keep all these pieces, called \emph{packets}, with
compatible product collars. Each label has two open \emph{forced cores},
one of each sign, on which that sign must be retained. A label is
\emph{decided} at a point when the point lies in one of its cores.
On an \(m\)-dimensional packet, positive
scalar curvature is required only where at least \(m-1\) remaining
labels are outside their forced cores. Minimizing walls may be
singular in high dimensions. Near the dangerous singular region,
additional pressure profiles make the remaining labels decided.
Hardy estimates, compactness across the bulk--cylinder seam, and a
maximum estimate for concentrating potentials pay for this shielding.
A guarded regularization and positive supersolutions then restore
smooth scalar-curvature control. The full finite-diagram arguments are
proved in \Cref{sec:packets,sec:cut,sec:local,sec:shield-global,sec:regularization}.
They include the stability and variation computations used here;
classical elliptic and geometric-measure-theory results are cited with
their hypotheses at their points of use.

Third, a countable list needs more than finiteness on compact subsets.
After a first radial cut, only finitely many labels and warp factors
are nonconstant on any \emph{whole} packet and on each finite incident
diagram needed by a step. All other labels are already forced; their
potentials form a summable list of spatial constants. Approximations
reproduce those constants exactly. This invariant permits the finite
packet arguments to be iterated with successively smaller positive
budgets, without a common positive floor for all labels. An intersection
of \(n-1\) distinct walls would produce a complete one-dimensional
packet. The positive scalar-curvature invariant on its finite warped
product yields a Riccati inequality that cannot hold on the complete
arclength line (\Cref{countable:dimension-one}), giving the required
wall-intersection bound.

Finally, the analytic cylinders are used only through finite
truncations in the topological reconstruction. Successive collar
compressions are supported over the corresponding original outer
balls. On one fixed copy of \(M\), the collapse maps and the old wall
statuses are therefore eventually unchanged near every point. This
local stabilization, together with a locally finite thickening of the
walls, gives continuity in the ordinary polyhedral topology. The
resulting constant is
\[
 C_n=4D+2,
\]
where the original scale \(D\) is chosen in terms of \(n\) before any
metric-dependent analytic estimates. The construction does not aim to optimize this constant.

The proof follows two separate inputs. Sections~\ref{pressure:section}--\ref{sec:rounds}
construct the geometric stretching and its scalar reserve.
Section~\ref{sec:finite-cuts} states the finite-label cutting procedure,
whose complete proof occupies Appendices~\ref{sec:packets}--\ref{sec:regularization}.
Section~\ref{sec:countable} proves its countable-label adaptation, and
Section~\ref{sec:reconstruction} constructs the width map.
Section~\ref{sec:closed-covers} verifies the uniform simplex metric and
proves the closed-cover consequence in dimensions two and three.

\begin{convention}\label{conv:signs}
The Laplacian is \(\Delta=\Div\grad\), and round spheres have positive
scalar curvature. For a hypersurface with unit normal \(\nu\), its
second fundamental form is
\(B(v,w)=\langle\nabla_v\nu,w\rangle\), and its mean curvature is
\(\tr B\). Integration and divergence use Riemannian densities, so no
orientation is needed. Distances in the final fiber bound are always
measured in the original metric \(g\).
\end{convention}
